\section{Задачи}
\begin{exersize}
	Исследовать функции $\sin{x}$, $\cos{x}$, $\ctg{x}$, $\arcsin{x}$, $\arccos{x}$, $\arctg{x}$, $\arcctg{x}$ по плану:
	\begin{enumerate}[1)]
		\item найти область определения; если область определения бесконечна, найти
			пределы на бесконечности или указать, что функция не имеет пределов на
			бесконечности;
		\item найти промежутки непрерывности;
		\item найти точки разрыва;
		\item найти односторонние пределы в точках разрыва;
		\item выяснить, является ли функция чётной, нечётной или ни чётной, ни
			нечётной;
		\item найти асимптоты;
		\item построить график.
	\end{enumerate}
\end{exersize}

\begin{exersize}
	Исследовать функции $x \mapsto 1/\sqrt{x^3}$, $x \mapsto 1/\sqrt[5]{x^4}$, $x
	\mapsto \sqrt[3]{x-2}$ по тому же плану.
\end{exersize}

\begin{exersize}
	Начертить эскиз графика функции $f(x)$ и найти предел $\lim\limits_{\baseB} f(x)$:
	\begin{align*}
		f(x) &= \tg \left( x + \frac{\pi}{3} \right), &
		\baseB &= \left( x \to \frac{\pi}{6} + 0 \right), &
		\baseB &= \left( x \to \frac{\pi}{6} - 0 \right);
		\\
		f(x) &= \log_{\frac{1}{3}} (x-2), &
		\baseB &= (x \to 2+0), &
		\baseB &= (x \to +\infty);
		\\
		f(x) &= \log_{\frac{1}{3}} (2-x), &
		\baseB &= (x \to 2-0), &
		\baseB &= (x \to -\infty);
		\\
		f(x) &= \left( \frac{1}{2} \right)^{3-x}, &
		\baseB &= (x \to -\infty), &
		\baseB &= (x \to +\infty);
		\\
		f(x) &= 4^{\tg{x}}, &
		\baseB &= \left( x \to \frac{\pi}{2} + 0 \right);
		\\
		f(x) &= \left( \frac{1}{3} \right)^{\frac{1}{2-x}}, &
		\baseB &= (x \to 2 - 0);
		\\
		f(x) &= \ln{\tg{x}}, &
		\baseB &= (x \to 0+0).
	\end{align*}
\end{exersize}

\begin{solution}
	Рассмотрим функцию~$x \mapsto \sin{x}$.
	\begin{enumerate}[1.]
		\item Область определения: $D(\sin{x}) = \RR$.
		\item Функция $f(x) = \sin{x}$ непрерывна на~$\RR$, то есть для любой
			точки~$x_0 \in \RR$
			\[
				\lim_{x \to x_0} \sin{x} = \sin{x_0}.
			\]
			Покажем это. Воспользуемся определением предела. Пусть $a_0 :=
			\sin{x_0}$, $\epsilon > 0$, $U_{\epsilon}(a_0)$ ---
			$\epsilon$-окрестность точки $a_0$, а $\Vo_{\delta}(x_0)$ --- проколотая
			$\delta$-окрестность точки~$x_0$.  Пусть $x \in \Vo_\delta(x_0)$;
			покажем, что по заданному~$\epsilon$ можно найти число~$\delta$ такое,
			что $\sin \bigl( \Vo_\delta(x_0) \bigr) \subset U_\epsilon (a_0)$. Так
			как максимальный по модулю тангенс угла наклона касательной к графику
			функции~$\sin{x_0}$ равен~$1$, то достаточно взять $\delta(\epsilon) =
			\epsilon$.

			Пределы на бесконечности не существуют. Чтобы это показать, перейдём от
			предела
			\[
				\lim_{x \to +\infty} {\sin{x}}
			\]
			к пределу
			\[
				\lim_{t \to 0 + 0}{\sin{\frac{1}{t}}}
			\]
			(мы положили $t := 1/x$).
			Пусть $a = \lim\limits_{t \to 0 + 0} \sin(1/t)$. Пусть $U_\epsilon(a)
			\cap [-1, 1]$ --- $\epsilon$-окрестность предельной точки в~$[-1, 1]$.

		\item Так как функция~$f$ непрерывна, то точек разрыва нет.

		\item Точек разрыва нет.

		\item Функция $x \mapsto \sin{x}$ --- нечётная функция, так как $\sin(-x) =
			-\sin{x}$.

		\item Так как предела при $x \to \pm \infty$ нет и точек разрыва нет, то
			асимптот нет.

		\item Чтобы построить график функции~$f(x) = \sin{x}$, определим:
			\begin{enumerate}[1)]
				\item нули функции~$f$ --- то есть точки, в которых~$f(x) = 0$;
				\item нули производной $f'$ функции~$f$ --- то есть точки, в которых
					$f'(x) = 0$;
				\item промежутки знакопостоянства $f$ (выше оси~$Ox$ или ниже
					расположен график?);
				\item промежутки знакопостоянства $f'$ (график едет вверх или вниз?).
			\end{enumerate}

			Поплыли.
			\begin{description}
				\item[Нули функции~$\bm{f}$.] 
					Функция $\sin{x}$ имеет нули в точках $\pi \ZZ = \set{\pi m}{m \in
					\ZZ}$; при переходе через каждый из нулей она меняет знак, так как
					каждый нуль --- простой%
					\footnote{Не является одновременно нулём производной.}%
					.

					Подставим значение $x = \pi/2$ из внутренности отрезка $[0, \pi]$ в
					функцию $f(x) = \sin{x}$. Получим, что $f(\pi/2) = \sin(\pi/2) = 1 >
					0$, и значит, что на всём интервале $(0, \pi)$ функция $\sin{x}$
					строго больше нуля.

				\item[Нули функции~$\bm{f'}$.]

					Так как $f'(x) = (\sin{x})' = \cos{x}$, а нули косинуса --- это
					\[
						\frac{\pi}{2} + \pi\ZZ :=
						\set{\frac{\pi}{2} + \pi m}{m \in \ZZ}
						.
					\]
			\end{description}
	\end{enumerate}
\end{solution}
